\otazka{4. Právní úprava kyberbezpečnosti v ČR a EU}{Základní instituty, principy, povinné orgány, systém zajištění kybernetické bezpečnosti.}{BVV03K}

\subsection*{4. Právní úprava kyberbezpečnosti -- stručněji}
Tuhle otázku je potřeba říct prakticky: kdo spadá do regulace, podle čeho se to pozná, jaký režim povinností má, co musí dělat, komu hlásí incidenty a jakou roli má NÚKIB.

\textbf{1. Co právo kyberbezpečnosti řeší}
\begin{itemize}
  \item Kyberbezpečnostní regulace neřeší každý počítač na světě. Soustředí se na regulované služby, významné organizace, základní bezpečnostní opatření, hlášení incidentů, dohled a spolupráci.
  \item Je potřeba odlišit kyberbezpečnost, kyberobranu, kyberkriminalitu a ochranu osobních údajů. Mohou se potkat u jednoho incidentu, ale mají jiné právní režimy.
\end{itemize}

\textbf{2. EU rámec}
\begin{itemize}
  \item NIS2 je hlavní směrnice pro kyberbezpečnost významných a základních subjektů. Stojí na řízení rizik, bezpečnostních opatřeních, hlášení incidentů, odpovědnosti vedení a dohledu.
  \item CER řeší odolnost kritických subjektů all-hazards způsobem. DORA řeší digitální provozní odolnost finančního sektoru. Cyber Resilience Act míří na bezpečnost produktů s digitálními prvky.
  \item Cyber Solidarity Act vytváří evropský systém kybernetického varování, cyber hubs, mechanismus kybernetické nouze, evropskou kybernetickou rezervu a přezkum incidentů.
\end{itemize}

\textbf{3. Český rámec}
\begin{itemize}
  \item Základem je zákon o kybernetické bezpečnosti č. 264/2025 Sb. a prováděcí vyhlášky. Praktické jádro je sebeidentifikace: subjekt má sám posoudit, zda poskytuje regulovanou službu.
  \item Nižší režim typicky dopadá na méně významné regulované služby, vyšší režim na významnější služby s přísnějšími opatřeními, audity a řízením bezpečnosti.
  \item Povinnosti zahrnují kontaktní údaje, bezpečnostní role, řízení rizik, evidenci aktiv, bezpečnostní opatření, řízení dodavatelů, školení, dokumentaci a hlášení incidentů.
\end{itemize}

\textbf{4. NÚKIB a nástroje}
\begin{itemize}
  \item NÚKIB je hlavní civilní správní autorita pro kybernetickou bezpečnost. Vykonává regulaci, dohled, metodickou podporu, koordinaci a provozuje GovCERT.CZ.
  \item Zákon rozlišuje běžná bezpečnostní opatření subjektů, nápravná opatření při dohledu a protiopatření NÚKIB: výstrahu, varování a reaktivní protiopatření.
  \item Stav kybernetického nebezpečí vyhlašuje NÚKIB jako úřad na své úřední desce, ne přímo ředitel osobně. Je na nezbytně nutnou dobu, nejvýše 30 dnů, s možností prodloužení v zákonných mezích do 60 dnů.
\end{itemize}

\subsection*{Kostra odpovědi}
\begin{enumerate}
  \item Odlišit kybernetickou bezpečnost, kyberobranu, kyberkriminalitu a ochranu osobních údajů.
  \item Vyložit EU rámec: NIS2, CER, DORA, CRA a související předpisy.
  \item Vyložit český rámec: zákon o kybernetické bezpečnosti a prováděcí vyhlášky.
  \item Popsat regulované služby, vyšší a nižší režim, sebeidentifikaci a hlášení incidentů.
  \item Popsat orgány a systém zajištění kybernetické bezpečnosti.
\end{enumerate}

\subsection*{Základní rozlišení}
Kybernetická bezpečnost je prevence, řízení rizik, ochrana systémů a služeb, detekce incidentů, reakce a obnova. Kyberobrana je vojenská dimenze kyberprostoru a patří do obranné politiky. Kyberkriminalita je trestná činnost a řeší ji orgány činné v trestním řízení. Ochrana osobních údajů chrání práva subjektů údajů a řeší ji hlavně GDPR a ÚOOÚ.

Jeden incident může spadat do více režimů. Ransomware v nemocnici může být kybernetický bezpečnostní incident podle ZKB, trestný čin podle trestního zákoníku, porušení zabezpečení osobních údajů podle GDPR a krizová situace z hlediska zdravotnictví.

\subsection*{EU rámec}
\textbf{NIS2} je hlavní směrnice EU pro vysokou společnou úroveň kybernetické bezpečnosti. Rozšiřuje okruh regulovaných subjektů, posiluje řízení rizik, odpovědnost vedení, hlášení incidentů, dohled a sankce. V ČR je transponována novým zákonem o kybernetické bezpečnosti.

\textbf{CER} řeší odolnost kritických subjektů v širším smyslu: fyzická, organizační a provozní odolnost. Je důležitá tam, kde narušení služby ohrožuje společnost bez ohledu na to, zda příčinou je kyberútok, sabotáž, přírodní událost nebo technické selhání.

\textbf{DORA} je zvláštní úprava pro finanční sektor. Zaměřuje se na řízení ICT rizik, hlášení incidentů, testování digitální provozní odolnosti a rizika spojená s externími ICT poskytovateli. U zkoušky je dobré říct, že DORA je lex specialis pro finance, ale neznamená, že finanční subjekt nikdy neřeší jiné předpisy.

\textbf{CRA} zavádí požadavky na produkty s digitálními prvky. Posouvá kyberbezpečnost do fáze návrhu, vývoje a uvádění produktů na trh. Výrobci budou řešit bezpečnostní požadavky, správu zranitelností, aktualizace a hlášení aktivně zneužívaných zranitelností.

Vedle těchto předpisů je dobré zmínit i institucionální vrstvu EU. \textbf{ENISA} podporuje členské státy, certifikaci, metodiky a evropskou spolupráci. \textbf{CSIRT Network} propojuje národní CSIRT týmy. \textbf{EU-CyCLONe} podporuje koordinaci rozsáhlých kybernetických krizí. \textbf{CERT-EU} chrání instituce EU. U trestních věcí vstupuje \textbf{Europol/EC3} a Eurojust. EU tedy nedělá jen normy, ale buduje i praktické koordinační a krizové mechanismy.

\textbf{Cybersecurity Act} - nařízení 2019/881 - je důležitý proto, že dal ENISA stálý mandát jako Agentuře EU pro kybernetickou bezpečnost a vytvořil evropský rámec kyberbezpečnostní certifikace. Certifikace se týká ICT produktů, služeb a procesů; může být dobrovolná, ale jiné předpisy nebo zadávací podmínky z ní mohou udělat prakticky nutný standard. Zkouškově si pamatuj: NIS2 řeší povinnosti subjektů, Cybersecurity Act řeší ENISA a certifikační rámec.

\textbf{Cyber Solidarity Act} - nařízení 2025/38 - posiluje unijní schopnosti detekce, připravenosti a reakce na velké kybernetické hrozby a incidenty. Staví na třech pilířích: \textbf{Evropský systém kybernetického varování}, \textbf{mechanismus kybernetické nouze} a \textbf{mechanismus přezkumu kybernetických incidentů}. Evropský systém kybernetického varování tvoří \textbf{národní kybernetická centra} (National Cyber Hubs) a \textbf{přeshraniční kybernetická centra} (Cross-Border Cyber Hubs). Tato centra mohou využívat funkce a schopnosti podobné bezpečnostním operačním centrům, například pokročilou detekci, sdílení informací a analýzu hrozeb, ale právní předpis je označuje jako cyber hubs, ne jako běžná SOC. Mechanismus kybernetické nouze zahrnuje mimo jiné \textbf{Evropskou kybernetickou rezervu}. Nejde o náhradu NIS2, ale o operační a krizovou nadstavbu pro společnou detekci, podporu a poučení z rozsáhlých incidentů.

\subsection*{Český zákon o kybernetické bezpečnosti}
Zákon o kybernetické bezpečnosti (č. 264/2025 Sb.) je česká transpozice NIS2. Neptá se obecně, jestli má firma "nějaké IT", ale jestli \textbf{poskytuje regulovanou službu}. Praktický postup je:
\begin{enumerate}
  \item Najdu službu ve \textbf{vyhlášce o regulovaných službách}. Typická odvětví jsou veřejná správa, energetika, doprava, zdravotnictví, digitální infrastruktura, finanční trh, voda, potraviny, chemie, výzkum a další.
  \item Ověřím \textbf{velikost a význam}. Typicky platí: \textbf{střední podnik} ve vybraném odvětví spadne do nižšího režimu, \textbf{velký podnik} nebo služba s vysokým dopadem do vyššího režimu. Vždy ale rozhoduje konkrétní služba a kritéria ve vyhlášce.
  \item Pokud podmínky sedí, subjekt službu \textbf{ohlásí NÚKIB}. Nečeká pasivně, až ho stát najde.
  \item Po registraci doplní kontaktní a další údaje, vymezí aktiva související s regulovanou službou a zavede bezpečnostní opatření podle svého režimu.
\end{enumerate}

\priklad{Střední výrobce v regulovaném výrobním odvětví si ověří službu ve vyhlášce a typicky řeší nižší režim. Velká nemocnice, významný provozovatel energetiky nebo zásadní digitální infrastruktura typicky řeší vyšší režim. Neříkej ale automaticky "každá firma od 50 lidí je NIS2"; vždy musí jít o regulovanou službu a splněná kritéria.}

\subsection*{Nižší a vyšší režim prakticky}
\begin{longtable}{p{0.18\linewidth}p{0.34\linewidth}p{0.40\linewidth}}
\toprule
\textbf{Režim} & \textbf{Kdo typicky} & \textbf{Co musí prakticky zavést}\\
\midrule
\textbf{Nižší povinnosti} &
Typicky střední podnik v regulovaném odvětví nebo méně významný poskytovatel regulované služby. &
Určí osobu pověřenou kybernetickou bezpečností, proškolí vedení a uživatele, sepíše přehled opatření, eviduje aktiva, vede jednoduché hodnocení rizik, nastaví pravidla přístupů a identit, řeší logování/detekci událostí, incidentní postup, kontinuitu a obnovu, základní síťovou a aplikační bezpečnost a používání kryptografie.\\
\addlinespace
\textbf{Vyšší povinnosti} &
Typicky velký nebo výrazně významný poskytovatel: podle velikosti, počtu uživatelů, geografického rozsahu, dopadu na odvětví, dopadu na jinou regulovanou službu nebo rizikovosti provozu. Do vyššího režimu spadne i služba registrovaná z moci úřední pro zvláštní význam. &
Zavede plnější ISMS: bezpečnostní politiku, role, řízení aktiv a rizik, řízení dodavatelů, HR bezpečnost, řízení změn, bezpečný vývoj a nákupy, řízení přístupů, zvládání incidentů, kontinuitu, audit. Technicky řeší fyzickou bezpečnost, segmentaci a bezpečnost sítí, správu identit, oprávnění, detekci, logování, vyhodnocování událostí, aplikační bezpečnost, kryptografii, dostupnost služby a OT/řídicí systémy, pokud je má.\\
\bottomrule
\end{longtable}

\textbf{Střední podnik} si pro zkoušku představ jako organizaci, která už není malá: zjednodušeně má méně než 250 zaměstnanců a zároveň roční obrat nejvýše 50 mil. EUR nebo bilanční sumu nejvýše 43 mil. EUR. \textbf{Velký podnik} už do této definice nespadá, typicky má 250 a více zaměstnanců nebo překročí finanční hranice. ZKB používá evropskou logiku velikosti podniku, ale u některých služeb a subjektů jsou zvláštní pravidla; proto je správná odpověď vždy: \emph{nejdřív služba a vyhláška, potom velikost a význam}.

\subsection*{Konkrétní opatření, která říct u zkoušky}
U \textbf{nižšího režimu} řekni hlavně toto:
\begin{itemize}
  \item vedení určí odpovědnou osobu pro kyberbezpečnost a pravidelně se seznamuje se stavem bezpečnosti,
  \item organizace ví, jaké služby, data, systémy, účty, dodavatelé a zařízení regulovanou službu drží nad vodou,
  \item existuje seznam hlavních rizik: ransomware, výpadek dodavatele, kompromitace admin účtu, ztráta záloh, zranitelná aplikace,
  \item přístupy se dávají podle potřeby, administrátorské účty jsou oddělené a důležité účty se silně ověřují,
  \item události se zaznamenávají a někdo je vyhodnocuje; incident má předem daný postup, kontakty a eskalační cestu,
  \item existují zálohy, pořadí obnovy služeb a nacvičený postup, co dělat při výpadku,
  \item zaměstnanci a administrátoři jsou školeni; dodavatelské smlouvy obsahují bezpečnostní požadavky.
\end{itemize}

U \textbf{vyššího režimu} přidej navíc:
\begin{itemize}
  \item plnohodnotný systém řízení bezpečnosti informací: politika, dokumentace, cíle, pravidelné přezkoumání,
  \item bezpečnostní role: manažer kybernetické bezpečnosti, architekt kybernetické bezpečnosti, garant aktiva a auditor kybernetické bezpečnosti,
  \item formální hodnocení rizik a plán zvládání rizik s termíny, odpovědnostmi a kontrolou účinnosti opatření,
  \item řízení významných dodavatelů: evidence, hodnocení rizik, bezpečnostní požadavky ve smlouvách, právo kontroly, řešení subdodavatelů,
  \item audit kybernetické bezpečnosti a pravidelné informování vedení,
  \item technické pokrytí od fyzické bezpečnosti přes sítě, IAM, logování, detekci a aplikační bezpečnost až po dostupnost služby a zabezpečení OT/ICS.
\end{itemize}

\subsection*{NÚKIB, protiopatření a dodavatelé}
NÚKIB je hlavní regulátor. V praxi vede registraci regulovaných služeb, kontroluje plnění povinností, přijímá hlášení ve vyšším režimu, vydává metodiky a může ukládat opatření. Vedle běžných bezpečnostních opatření existují přesné druhy protiopatření: \textbf{výstraha}, \textbf{varování} a \textbf{reaktivní protiopatření}. Výstraha upozorňuje na nebezpečí, varování na hrozbu, kterou mají dotčené subjekty zohlednit v řízení rizik, a reaktivní protiopatření už ukládá konkrétní technický nebo organizační krok.

Dodavatelský řetězec vysvětli jednoduše: regulovaný subjekt nemůže říct "to je problém našeho dodavatele". Pokud dodavatel spravuje nemocniční systém, cloud, síť, vzdálený přístup nebo OT technologii, jeho selhání může shodit regulovanou službu. Proto se řeší bezpečnostní požadavky ve smlouvě, přístupy dodavatele, subdodavatelé, auditovatelnost, reakce na incident a ukončení spolupráce.

\subsection*{Hlášení incidentů prakticky}
Vyšší a nižší režim se liší i hlášením. \textbf{Vyšší režim} hlásí širší okruh incidentů postihujících regulovanou službu, pokud mají původ v kyberprostoru a nelze rychle vyloučit úmyslné zavinění. \textbf{Nižší režim} hlásí užší okruh: incident musí mít významný dopad na poskytování regulované služby. Zkouškově stačí říct logiku: nejdřív rychlá informace, potom doplnění detailů, průběžné aktualizace podle potřeby a závěrečná zpráva.

Hlášení podle ZKB neznamená, že odpadá GDPR nebo trestní rovina. Ransomware v nemocnici může být současně kybernetický incident, trestný čin, porušení zabezpečení osobních údajů a problém kontinuity zdravotní péče.

\textbf{Stav kybernetického nebezpečí} je mimořádný režim pro závažné ohrožení kybernetické bezpečnosti; vyhlašuje jej Národní úřad pro kybernetickou a informační bezpečnost, tedy Úřad ve smyslu zákona č. 264/2025 Sb., na své úřední desce, a to na nezbytně nutnou dobu, nejvýše na 30 dnů. Lze jej prodloužit tak, aby celková doba nepřekročila 60 dnů.

\pravobox{Principy regulace: \textbf{technologická neutralita} znamená, že právo nemá být navázané na jednu konkrétní technologii. \textbf{Proporcionalita} znamená, že povinnosti a zásahy musí odpovídat riziku a dopadu. \textbf{Risk-based approach} znamená, že organizace má řídit skutečná rizika, ne jen formálně plnit seznam. \textbf{Odpovědnost vedení} znamená, že kyberbezpečnost není jen problém IT oddělení. \textbf{Autonomie regulovaného subjektu} znamená, že stát stanoví rámec, ale subjekt si musí prakticky nastavit vlastní řízení bezpečnosti.}

\subsection*{Praktické situace použití předpisů}
\priklad{Modelový příklad - nemocnice a ransomware: nemocnice zjistí ransomware a možné odcizení zdravotnické dokumentace. Pokud poskytuje regulovanou službu, řeší podle ZKB kybernetický bezpečnostní incident, hlášení NÚKIB a splnění bezpečnostních opatření; vedení nemocnice nemůže tvrdit, že jde jen o problém IT. Současně běží GDPR: zdravotní údaje jsou zvláštní kategorie, správce posuzuje porušení zabezpečení, oznamuje ÚOOÚ při riziku a informuje pacienty při vysokém riziku. Policie řeší trestní rovinu, nemocnice kontinuitu péče, zřizovatel kapacity a komunikaci. Pokud jde zároveň o kritickou infrastrukturu nebo kritický subjekt, vstupuje i CER/ZKI logika odolnosti. Ke zkoušce: jeden incident spouští více režimů vedle sebe, žádný z nich automaticky "nepřebije" všechny ostatní.}

\priklad{Modelový příklad - banka a cloudový poskytovatel: banka ztratí dostupnost klíčového systému kvůli selhání externího ICT poskytovatele. Základní rámec pro finanční sektor je DORA: řízení ICT rizik, evidence a hlášení závažného ICT incidentu, řízení rizik třetích stran, smluvní požadavky, plány kontinuity a testování obnovy. Pokud incident zároveň obsahuje osobní údaje, běží GDPR. Pokud banka nebo dodavatel spadá i do obecné kybernetické regulace, řeší se souběh podle konkrétní role a služby; není přesné říct, že DORA vždy úplně ruší všechny ostatní předpisy. Ke zkoušce: DORA je lex specialis pro digitální provozní odolnost financí, ale souběžné režimy se posuzují podle účelu a dopadu.}

\priklad{Modelový příklad - výrobce připojeného routeru a CRA: výrobce zjistí aktivně zneužívanou zranitelnost ve firmware routeru. Cyber Resilience Act posouvá odpovědnost do celého životního cyklu produktu: bezpečný návrh, dokumentace, správa zranitelností, bezpečnostní aktualizace, informování uživatelů a posouzení shody. K 18. 6. 2026 je CRA v platnosti, ale plná sada hlavních povinností pro produkty s digitálními prvky se použije až od 11. 12. 2027; zvláštní hlášení aktivně zneužívaných zranitelností a závažných incidentů začne dříve, od 11. 9. 2026. Ke zkoušce: k 18. 6. 2026 nejde tvrdit, že všechny povinnosti CRA už prakticky běží.}

\priklad{Modelový příklad - Cyber Solidarity Act: rozsáhlý přeshraniční útok zasáhne poskytovatele služeb ve více členských státech. Národní a přeshraniční kybernetická centra v Evropském systému kybernetického varování sdílejí detekční poznatky a korelují hrozby; při závažném incidentu může mechanismus kybernetické nouze pomoci s podporou, například přes Evropskou kybernetickou rezervu, a po incidentu může následovat přezkum kybernetického incidentu. Rezerva ale není náhradou odpovědnosti provozovatele: ten musí dál řídit rizika, obnovu, hlášení a komunikaci. Ke zkoušce: EU mechanismus pomáhá koordinovat a posílit kapacity, nenahrazuje lokální povinnosti subjektu.}

\subsection*{Orgány a systém}
\begin{itemize}
  \item \textbf{NÚKIB} - gestor kybernetické bezpečnosti, dohled, varování, opatření, metodiky, GovCERT.CZ.
  \item \textbf{GovCERT.CZ} - vládní tým pro řešení incidentů.
  \item \textbf{CSIRT.CZ} - národní CSIRT, koordinace incidentů v širším ekosystému.
  \item \textbf{Regulované subjekty} - nesou primární odpovědnost za bezpečnost svých služeb.
  \item \textbf{Vrcholové vedení} - schvaluje a nese odpovědnost za bezpečnostní opatření; NIS2 posiluje osobní odpovědnost managementu.
  \item \textbf{Policie, zpravodajské služby, armáda, ÚOOÚ, ČTÚ, DIA, NBÚ a další} - vstupují podle povahy incidentu a právního režimu.
\end{itemize}

Prakticky se dá systém popsat jako vícevrstvý. NÚKIB a GovCERT řeší civilní kyberbezpečnost a regulaci. CSIRT.CZ podporuje širší technickou komunitu. Policie a státní zastupitelství řeší trestní rovinu. Zpravodajské služby pomáhají s hrozbami, špionáží a atribucí. Armáda a Vojenské zpravodajství řeší obrannou dimenzi. ÚOOÚ vstupuje při osobních údajích, ČTÚ u elektronických komunikací, DIA u digitálních služeb státu a NBÚ u utajovaných informací.

\priklad{Výrobní podnik nejdřív neřeší "jsme NIS2?", ale konkrétně: poskytujeme regulovanou službu podle české vyhlášky? V jakém odvětví? Splňujeme velikostní nebo významová kritéria? Až potom určí režim povinností.}

\priklad{Nemocnice ve vyšším režimu musí uvažovat nejen antivir a firewall, ale i odpovědnost vedení, řízení dodavatelů nemocničního systému, zálohy oddělené od produkce, krizové postupy pro výpadek IT, školení personálu, hlášení incidentu NÚKIB a souběžně GDPR režim, pokud uniknou údaje pacientů.}

\pozor{Neříkej, že NIS2 přímo ukládá české firmě všechny povinnosti. Jako směrnice potřebuje transpozici; v ČR je klíčový zákon o kybernetické bezpečnosti (č. 264/2025 Sb.) a prováděcí vyhlášky.}

\zdroje{Směrnice 2022/2555 (NIS2); zákon o kybernetické bezpečnosti (č. 264/2025 Sb.); vyhlášky k regulovaným službám, vyšším povinnostem a nižším povinnostem; DORA 2022/2554; CRA 2024/2847.}

\newpage

